% Triangle Mutiny — Modal Adjacency
% Author: Borealis S. Hedling
% Version: 1.0

\section*{Modal Adjacency}

Modal adjacency defines which semantic modes may transition into one another
without leaving the stable manifold. Modal structure is discrete: each point
\( p \in M \) carries a modal label \( \mu(p) \in L \), where \( L \) is a finite
set of modes. Adjacency constraints prevent semantic mimicry and modal drift.

\subsection*{1. Modal Set}

Let the modal set be:
\[
L = \{ \ell_1, \ell_2, \ldots, \ell_n \}.
\]

Modes represent semantic neighborhoods. They do not encode meaning directly;
they encode structural position within the reasoning manifold.

\subsection*{2. Adjacency Relation}

Define the adjacency relation:
\[
E \subseteq L \times L.
\]

A pair \( (\ell_i, \ell_j) \in E \) indicates that a transition from mode
\( \ell_i \) to mode \( \ell_j \) is structurally valid.

Adjacency is not symmetric unless explicitly defined.

\subsection*{3. Valid Transition}

A transition \( p^- \rightarrow p^+ \) is valid when:
\[
(\mu(p^-), \mu(p^+)) \in E.
\]

Valid transitions preserve modal coherence and prevent semantic drift.

\subsection*{4. Invalid Transition}

A transition is invalid when:
\[
(\mu(p^-), \mu(p^+)) \notin E.
\]

Invalid transitions indicate modal instability, semantic mimicry, or collapse
surface proximity.

\subsection*{5. Modal Collapse Predicate}

Define modal collapse:
\[
C_{\mathrm{modal}}(p) =
\begin{cases}
\mathrm{true} & \text{if } (\mu(p^-), \mu(p^+)) \notin E, \\
\mathrm{false} & \text{otherwise}.
\end{cases}
\]

This predicate integrates directly with the unified collapse predicate defined
in collapse-geometry.tex.

\subsection*{6. Modal Neighborhood}

For each mode \( \ell \in L \), define its neighborhood:
\[
N(\ell) = \{ \ell' \in L \mid (\ell, \ell') \in E \}.
\]

Modal neighborhoods determine the local structure of the modal graph.

\subsection*{7. Modal Graph}

The modal graph is the directed graph:
\[
G = (L, E).
\]

The modal graph encodes all valid transitions. Cycles in \( G \) represent
modal loops; acyclic regions represent hierarchical modal structure.

\subsection*{8. Modal Stability Condition}

A reasoning trajectory \( \gamma = (p_1, p_2, \ldots, p_k) \) is modal-stable
when:
\[
(\mu(p_i), \mu(p_{i+1})) \in E
\quad \text{for all } i.
\]

Violation indicates modal drift.

\subsection*{9. Modal Drift Magnitude}

Define modal drift magnitude:
\[
\Delta_{\mathrm{modal}}(p) =
\begin{cases}
0 & \text{if transition is valid}, \\
1 & \text{if transition is invalid}.
\end{cases}
\]

This primitive feeds into holonomy drift tensors.

\section*{Summary}

This file defines modal labels, adjacency relations, valid and invalid
transitions, modal neighborhoods, modal graphs, and modal collapse predicates.
Modal adjacency is a core structural invariant of Triangle Mutiny and integrates
directly with collapse geometry and stability envelopes.
